% Source: Weinberg GR, physical PDF pp. 428--430
%         (printed pp. 407--409).
% Coverage: Chapter 14 epigraph and opening overview before Section 14.1.
% Figures/tables/footnotes: epigraph; linked reference markers 1--9; no
%                          figures, tables, or footnotes.
% Status: source-reviewed and compile-clean; visually proofed.
% Uncertainties: none.

\begin{flushleft}
\begin{minipage}{0.35\textwidth}
\raggedright
\itshape
``We wonder, Oh, we wonder,\\
what on earth the world\\
may be?''

\medskip
\upshape
\textit{W. S. Gilbert, The Mikado}
\end{minipage}
\end{flushleft}

Modern science began with the discovery that the earth is not at the center of
the universe. Antianthropocentrism has become incorporated into the scientific
mentality, and no one now would seriously suggest that the earth, or the solar
system, or our galaxy, or our local group of galaxies, occupies any specially
favored position in the cosmos. Rather, our intuition now runs in precisely the
opposite direction. A large portion of modern cosmological theory is built on
the \emph{Cosmological Principle}, the hypothesis that all positions in the
universe are essentially equivalent. Of course, the homogeneity of the universe
has to be understood in the same sense as the homogeneity of a gas: It does not
apply to the universe in detail, but only to a ``smeared-out'' universe averaged
over cells of diameter \(10^8\) to \(10^9\) light years, which are large enough
to include many clusters of galaxies. Also, it appears that the universe is
spherically symmetric about us, so included in the Cosmological Principle is
the assumption that the ``smeared'' universe is isotropic about every point.

The question still remains whether the universe is spherically symmetric and
homogeneous at all times, or merely over some temporary present phase of its
history. There has been an interesting suggestion, to be discussed in
Section~15.11, that the universe may have been highly anisotropic during some
dense early phase, but that the anisotropies have since been largely smoothed
out through the action of neutrino viscosity and other dissipative effects.
However, even in such theories, the universe has been highly isotropic and
homogeneous over all of that part of its history which is directly accessible
to astronomical observation.

This chapter will outline and apply a mathematical framework for the
description of the universe, based entirely on the Cosmological Principle, and
on those parts of general relativity that follow directly from the Principle of
Equivalence. (This includes Chapters~2 to~6 and Chapter~13.) I shall first show
that the Cosmological Principle allows the specification of the cosmic metric
entirely in terms of a scale factor \(a(t)\) and a trichotomic constant \(k\)
[as in Eq.~\eqref{eq:13.5.32}], and we shall then see how astronomical
observations can be interpreted as measurements of \(a(t)\) and \(k\).

This inherently kinematic approach, pioneered in the 1930s by H. P.
Robertson\hyperref[ch14-ref-1]{[1]} and A. G.
Walker,\hyperref[ch14-ref-2]{[2]} is incomplete, in that it does not provide an
\emph{a priori} prediction of the function \(a(t)\). To calculate \(a(t)\) we
need to make some assumption about the material content of the universe, and
then derive the Robertson--Walker metric as a solution of the Einstein field
equations, as first done by Alexandre
Friedmann\hyperref[ch14-ref-3]{[3]} in 1922. Our discussion of the contents of
the universe, and the use of the Einstein field equations, will be postponed
until the next chapter, on cosmology.

Why make this distinction between cosmography and cosmology? The reason is
simply that we do not know the equation of state of the matter and radiation of
the universe throughout its history, and even if we did, we could not be sure
that the Einstein equations really hold over cosmic times and distances. A
modification of the field equations or the equation of state, such as the
introduction of a Brans--Dicke field, a cosmological constant, or a large
population of neutrinos or gravitons, would affect the function \(a(t)\) and
invalidate the simplest Friedmann solution, but it would not require us to make
any change in the descriptive framework assembled in this chapter.

There remains the possibility that the universe is not homogeneous and
isotropic after all. It might be homogeneous but not isotropic, as in the model
of K. G\"odel.\hyperref[ch14-ref-4]{[4]} However, the cosmic microwave
radiation discussed in Chapter~15 appears to be highly isotropic. (The universe
cannot be isotropic about every point without also being homogeneous, as shown
in the last chapter.) A more radical notion is that there is no ``smeared''
universe at all, but only clusters of galaxies, and clusters of clusters, and
clusters of clusters of clusters, and so on, as in the hierarchical model
proposed in 1908 by C. V. I.
Charlier.\hyperref[ch14-ref-5]{[5]} Empirical arguments for such
superclustering have been offered by G. de
Vaucouleurs,\hyperref[ch14-ref-6]{[6]} but the work of F.
Zwicky,\hyperref[ch14-ref-7]{[7]} G. O.
Abell,\hyperref[ch14-ref-8]{[8]} and J. H.
Oort\hyperref[ch14-ref-9]{[9]} indicates that the hierarchy stops at clusters
of galaxies or at most clusters of clusters of galaxies, and shows no evidence
of inhomogeneities of larger scale.

The real reason, though, for our adherence here to the Cosmological Principle
is not that it is surely correct, but rather that it allows us to make use of
the extremely limited data provided to cosmology by observational astronomy.
If we make any weaker assumptions, as in the anisotropic or hierarchical
models, then the metric would contain so many undetermined functions (whether
or not we use the field equations) that the data would be hopelessly inadequate
to determine the metric. On the other hand, by adopting the rather restrictive
mathematical framework described in this chapter, we have a real chance of
confronting theory with observation. If the data will not fit into this
framework, we shall be able to conclude that either the Cosmological Principle
or the Principle of Equivalence is wrong. Nothing could be more interesting.
